\section{A finite interchange map for the residual character at three}

This section compares the exceptional character $\chi_3$ with two
finite specializations of the harmonic and equianharmonic formulas of
Nakamura and Tsunogai.  An explicit homomorphism between finite matrix
groups places both comparisons in the same quaternion central sign.
Furusho's dyadic identities will identify the harmonic comparison
on full GT in Section~\ref{sec:harmonic-dyadic-comparison}. The
cubic comparison follows in Section~\ref{sec:normalized-cubic-identity}.

Write $G=\widehat{\mathrm{GT}}$ and let
$D=\ker\lambda\cap\bigcap_q\ker\rho_q^I$.  The pure harmonic and
cubic conjugators are denoted $\beta_g,A_g\in\widehat F(x,y)$,
respectively, where $x=u^2,y=v^2$ in the modular group.
We use
\[
 U=\begin{pmatrix}1&1\\0&1\end{pmatrix},\quad
 L=\begin{pmatrix}1&0\\-1&1\end{pmatrix},\quad
 S=ULU,\quad R=-UL.
\]
On $D$ the established signed normalization is $D_g=I$, and the
established finite frame targets give
\[
 q_3(\beta_g)^{-1}=(-I)^{w(g)},\qquad
 q_8(A_g)^{-1}=5^{\tau(g)}I,\qquad
 \chi_3(g)=w(g)+\tau(g).
\]
Here $q_n$ is the signed braid representation reduced modulo $n$.
All signs in this section are in $\mathbb F_2$.  In particular there is
no extra minus sign in the definition of the harmonic frame.

\begin{proposition}[Finite interchange map]\label{prop:finite-interchange}
Put $X=U^2\bmod8$ and
$P=\langle X,R\bmod8\rangle\subset\mathrm{SL}_2(\mathbb Z/8)$.
There is a surjection
\[
 \kappa:P\longrightarrow\mathrm{SL}_2(\mathbb F_3)
\]
such that
\[
 \kappa(X)=-S\bmod3,\qquad
 \kappa(R)=U\bmod3,\qquad
 \kappa(5I)=-I.
\]
The group $P$ has order $192$ and the kernel has order eight.
For every profinite commutator word $f\in\widehat F(x,y)'$,
\[
 \kappa\bigl(f(X,R)\bigr)=f(S,U)=f(U,S)^{-1}
\]
if $f$ satisfies the two-cycle relation.
\end{proposition}
\begin{proof}
Put $Y=L^2\bmod8$, $C=5I$ and $Z=-I$.  Direct matrix
multiplication gives
\[
 [X,Y]=C,\qquad X^4=Y^4=C^2=Z^2=I,
\]
and $X^2,Y^2,C,Z$ are central in $N=\langle X,Y,Z\rangle$.
Consequently every element of $N$ has the form
$X^aY^bC^eZ^d$ with $0\le a,b<4$ and $e,d\in\{0,1\}$.
These 64 forms are distinct.  Indeed, if one equals $I$, its two
off-diagonal entries give $2a=2b=0\bmod8$, so $a=b=0$;
then $5^e7^d=1\bmod8$ gives $e=d=0$.
This also proves that the displayed class-two presentation is a
presentation of $N$.

Conjugation by $R$ has the following action:
\[
 X\longmapsto Y,\qquad
 Y\longmapsto Z(XY)^{-1},\qquad C\longmapsto C,\qquad
 Z\longmapsto Z.
\]
Since $R^3=I$ and its reduction modulo two has order three,
whereas every element of $N$ is one modulo two, one has
$P=N\rtimes\langle R\rangle$, of order $64\cdot3$.
Here $Y=RXR^{-1}$ and
$Z=(RYR^{-1})XY$, so this semidirect product is generated by
$X,R$ as asserted.

In $\mathrm{SL}_2(\mathbb F_3)$ put
\[
 i=-S=\begin{pmatrix}0&2\\1&0\end{pmatrix},\qquad
 t=U,\qquad
 j=tit^{-1}=\begin{pmatrix}1&1\\1&2\end{pmatrix}.
\]
Then $i^2=j^2=[i,j]=-I$, $t^3=I$, and
$tjt^{-1}=ij$.  Thus $i,j$ generate the quaternion group and
$t$ acts on its three noncentral cyclic subgroups by a three-cycle.
The assignment
\[
 X\mapsto i,\quad Y\mapsto j,\quad C\mapsto-I,\quad Z\mapsto-I
\]
respects the presentation of $N$.  It also respects conjugation by
$R$, because
$(-I)(ij)^{-1}=ij$.  It therefore extends with $R\mapsto t$
to a homomorphism $P\to Q_8\rtimes C_3$
$=\mathrm{SL}_2(\mathbb F_3)$.  It is surjective, and its kernel
has order $192/24=8$.

Evaluation of a profinite word commutes with a homomorphism to a
finite group.  Multiplying either argument of a commutator word
by a central element does not change its evaluation, since the
two exponent sums are zero.  Hence
$\kappa(f(X,R))=f(-S,U)=f(S,U)$, and the two-cycle relation gives
the final equality.
\end{proof}

\begin{lemma}[A common finite evaluation]\label{lem:common-sign}
For every $g\in D$ the matrix $f_g(U,S)\bmod3$ is central.
Consequently
\[
 f_g(U,S)=(-I)^{\eta(g)}\quad\text{in }\mathrm{SL}_2(\mathbb F_3)
\]
defines a continuous homomorphism $\eta:D\to\mathbb F_2$.
\end{lemma}
\begin{proof}
Write $[f_g]=A_f(X,Y)[x,y]$ in the completed Alexander module.
The two-cycle relation gives $A_f(X,Y)=A_f(Y,X)$.
We recall precisely the cyclic-cover input already established for
full GT.  In the homology of the subgroup
$\ker(\widehat F(x,y)\to C_3)$ with $x\mapsto1,y\mapsto0$,
use the basis $x^3,y_j=x^jyx^{-j}$, $0\le j<3$.
The three $y_j$ coefficients of $f_g$ are
$(2b,-b,-b)$, where $b=-\rho_3^I(g)/2$.
Sending $y_j$ to $\omega^j$ gives
\[
 (\omega-1)A_f(\omega,1)=3b.
\]
On $D$, $b=0$.  Since $\omega-1$ is a unit in
$\mathbb Z_2[\omega]$, it follows that
$A_f(\omega,1)=A_f(1,\omega)=0$, also after reduction modulo two.

Now project $\mathrm{SL}_2(\mathbb F_3)$ to
$A_4=V\rtimes C_3$, where $V=Q_8/\{\pm I\}$ is identified
with the additive group of $\mathbb F_4$.
The image of $S$ belongs to $V$ and the image of $U$ acts on
$V$ by a primitive cube root $\omega$ or $\omega^{-1}$.
Conjugation by $S$ is trivial on $V$.  Since this quotient is
metabelian, the image of $f_g(S,U)$ is
$A_f(1,\omega)[S,U]$, with either choice of the root; it is zero.
Thus $f_g(S,U)$, and hence $f_g(U,S)$ by the two-cycle relation,
lies in $\{\pm I\}$.

For $g\in D$, the free-group automorphism
$x\mapsto x$, $y\mapsto f_g^{-1}yf_g$ induces the identity under
the evaluation $x\mapsto U,y\mapsto S$, because the evaluated
conjugating word is central.  Applying this observation to the
composition law $f_{gh}=f_g\phi_g(f_h)$ proves that the evaluated
signs multiply.  Continuity follows from evaluation in the finite
group $\mathrm{SL}_2(\mathbb F_3)$.
\end{proof}

\begin{proposition}[An exact pair of finite comparison gates]
\label{prop:pair-of-gates}
For $g\in D$ define
\[
 \epsilon_h(g)=w(g)+\eta(g),\qquad
 \epsilon_c(g)=\tau(g)+\eta(g).
\]
These are continuous homomorphisms on $D$, and
\[
 \boxed{\quad \chi_3=\epsilon_h+\epsilon_c.\quad}
\]
More explicitly, the two signs are the following finite defects:
\[
 \begin{split}
 (-I)^{\epsilon_h(g)}
   &=q_3(\beta_g)^{-1}f_g(U,S)^{-1},\\
 (-I)^{\epsilon_c(g)}
   &=\kappa\!\left(q_8(A_g)^{-1}f_g(X,R)^{-1}\right).
 \end{split}
\]
In particular the two finite comparisons
\begin{equation}
 q_3(\beta_g)^{-1}=f_g(U,S),\qquad
 \kappa(q_8(A_g)^{-1})=\kappa(f_g(X,R))
 \label{eq:finite-harmonic-cubic-comparisons}
\end{equation}
would imply $\chi_3(g)=0$.  They are sufficient conditions;
vanishing of $\chi_3$ requires only equality of their two defects.
Both defects vanish on $D\cap G_{\mathbb Q}$.
\end{proposition}
\begin{proof}
The established frame homomorphisms make $w,\tau$ homomorphisms
on $D$, and Lemma~\ref{lem:common-sign} does the same for $\eta$.
The first finite formula follows directly from their definitions.
For the second, Proposition~\ref{prop:finite-interchange} gives
\[
 \kappa(q_8(A_g)^{-1})=(-I)^{\tau(g)},\qquad
 \kappa(f_g(X,R))=f_g(U,S)^{-1}=(-I)^{\eta(g)}.
\]
This also proves the displayed sum and the sufficient conditions.

For the arithmetic assertion use exactly the scope of
\cite[Theorem~2, pp.~199--200]{NT2002}.  Its $(\mathrm{GF}_0)$ and
$(\mathrm{HF}_1)$ formulas are stated for Galois elements.
When $\lambda=1$ and $\rho_2=\rho_3=0$, they reduce to
\[
 \beta_g^{-1}=f_g(u,s),\qquad A_g^{-1}=f_g(u^2,r).
\]
Here its harmonic parameter is $\beta^{-1}$.  Its cubic
parameter is $A^{-1}$, since $\operatorname{Ad}(r)$ cyclically
permutes $x,y,(xy)^{-1}$ and
$f_g=\operatorname{Ad}(r)(A_g)A_g^{-1}$.
These identifications also follow from the uniqueness of the
pure conjugators.  The equalities lift without a central error
from the modular quotient to the signed braid group: both
sides have degree zero in the braid abelianization, whereas
the central generator has degree six.
Under the signed representation $q(r)=UL=-R$; the central
minus sign disappears in a commutator word.
Reduction modulo three and eight therefore gives
\eqref{eq:finite-harmonic-cubic-comparisons} on the arithmetic
subgroup.  No use of those two formulas on all of GT is made.
\end{proof}

\subsection{Two relative quadratic characters}

Continue with the same notation:
$G=\widehat{\mathrm{GT}}$, $H=G_{\mathbb Q}$,
$D=\ker\lambda\cap\bigcap_q\ker\rho_q^I$, and
$f_g(U,S)=(-I)^{\eta(g)}$ for $g\in D$.
The two frame signs are $w,\tau$, and
$\epsilon_h=w+\eta$, $\epsilon_c=\tau+\eta$.
We use the previously established facts $G=DH$, the invariance
of $D$, and the signed frame conjugation laws.  The latter make
$w$ and $\tau$ invariant under $G$: on $D$ the finite frames are
scalar matrices, which every conjugation fixes.

\begin{lemma}\label{lem:eta-quotient-invariant}
The kernel of the evaluation
\[
 e:\widehat F(x,y)\longrightarrow Q=\mathrm{SL}_2(\mathbb F_3),
 \qquad e(x)=U,\quad e(y)=S,
\]
is invariant under every gentle automorphism.  Consequently $\eta$
is invariant under conjugation by all of $G$.
\end{lemma}
\begin{proof}
We first verify the precise finite presentation:
\begin{equation}\label{eq:finite-quotient-presentation}
Q=\langle x,y\mid x^3=y^4=(xy)^3=1,\ [x,y^2]=1\rangle.
\end{equation}
In the presented group, $z=y^2$ is central and has order at most
two.  Modulo $z$ one has
$\langle x,y\mid x^3=y^2=(xy)^3=1\rangle$, of order at most
twelve.  For an explicit bound, put
$b_j=x^jyx^{-j}$ for $j=0,1,2$.  They are involutions, their
product in cyclic order is one by $(xy)^3=1$, and hence two
of them commute and generate a group of order at most four.
This group is normalized by $x$, which cyclically permutes the
$b_j$.  Thus the whole quotient has order at most $4\cdot3$.
The group in \eqref{eq:finite-quotient-presentation} has order at most 24.  The matrices $U,S$ satisfy
all four relators.  They generate $\mathrm{SL}_2(\mathbb F_3)$:
$-S$ and its conjugate by $U$ generate the quaternion subgroup,
and $U$ has order three and permutes its three noncentral cyclic
subgroups.  This gives 24 distinct matrices, proving \eqref{eq:finite-quotient-presentation}.

Let $g=(\lambda,f)$ be a gentle element and use the convention
\[
 \phi_g(x)=x^\lambda,\qquad
 \phi_g(y)=f^{-1}y^\lambda f.
\]
The peripheral three-cycle condition says that $\phi_g(xy)$
is conjugate to $(xy)^\lambda$.  Therefore the images under
$e\phi_g$ satisfy the first three relators of \eqref{eq:finite-quotient-presentation}.
Moreover
\[
 e\phi_g(y^2)=e(f)^{-1}(S^2)^\lambda e(f)=-I,
\]
since $S^2=-I$ is central and $\lambda$ is odd.
This proves the fourth relator as well.  There is consequently
an induced endomorphism $\alpha_g$ of $Q$ with
$e\phi_g=\alpha_ge$.  Applying the same argument to $g^{-1}$
shows that $\alpha_g$ is an automorphism.
In particular every $\alpha_g$ fixes the central involution
$-I$, directly also visible in the displayed formula.

Put $s(g)=e(f_g)$.  The composition convention is
$f_{gh}=f_g\phi_g(f_h)$, so
\[
 s(gh)=s(g)\alpha_g(s(h)).
\]
For $d\in D$, Lemma~\ref{lem:common-sign} proves
$s(d)\in\{\pm I\}$.  Since $\lambda_d=1$, the formulas on
the two generators then give $\alpha_d=\mathrm{id}_Q$.
For $g\in G$ and $d\in D$, it follows that
\[
 \begin{split}
 s(gdg^{-1})
 &=s(g)\alpha_g(s(d))\alpha_{gd}(s(g^{-1}))\\
 &=s(g)s(d)\alpha_g(s(g^{-1}))=s(d).
 \end{split}
\]
The last equality uses centrality of $s(d)$ and
$s(g)\alpha_g(s(g^{-1}))=1$.
This proves the conjugation invariance of $\eta$.
\end{proof}

\begin{proposition}\label{prop:two-relative-extensions}
Each of $\epsilon_h,\epsilon_c$ has a unique extension to a
continuous homomorphism $G\to\mathbb F_2$ that is zero on $H$.
Their sum is the established residual relative character $\chi_3$.
The harmonic extension is proved to vanish in
Section~\ref{sec:harmonic-dyadic-comparison}. The cubic extension
and $\chi_3$ vanish by Section~\ref{sec:normalized-cubic-identity}.
\end{proposition}
\begin{proof}
By Lemma~\ref{lem:eta-quotient-invariant} and the scalar frame
conjugation laws, both homomorphisms on $D$ are invariant under
$G$.  Proposition~\ref{prop:pair-of-gates}, using only the Galois
scope of the published $(\mathrm{GF}_0)$ and $(\mathrm{HF}_1)$
formulas, proves that both are zero on $D\cap H$.

For either homomorphism $\epsilon$, define
\[
 \widetilde\epsilon(dh)=\epsilon(d),\qquad d\in D,\ h\in H.
\]
If $dh=d'h'$, then $(d')^{-1}d\in D\cap H$, so this is
well-defined.  Conjugation invariance and normality of $D$ give
\[
 \widetilde\epsilon(d_1h_1d_2h_2)
 =\epsilon(d_1)+\epsilon(h_1d_2h_1^{-1})
 =\epsilon(d_1)+\epsilon(d_2).
\]
Thus it is a homomorphism, zero on $H$.
The continuous surjection $D\times H\to G$ is a quotient map,
because its source is compact and its target is Hausdorff.
Its pullback is the continuous function $(d,h)\mapsto\epsilon(d)$,
which proves continuity.  The decomposition $G=DH$ proves
uniqueness.  The two extensions sum to $w+\tau$ on $D$ and
are zero on $H$, hence their sum is $\chi_3$ by the same uniqueness.
\end{proof}

\subsection{A limitation of a low-degree peripheral calculation}

The following completely explicit test records why checking the
modular order-three relation through degree four does not remove
the central ambiguity.  It is not a counterexample to that relation
in the full profinite group and is not a GT counterexample.

Let $\Pi=\langle a,b\rangle$ be the discrete free group, with
automorphisms
\[
 \theta(a)=a^{-1},\quad\theta(b)=b^{-1},\qquad
 \alpha(a)=a^{-1}b,\quad\alpha(b)=a^{-1}.
\]
These are the actions of the order-two and order-three modular
generators on the free subgroup $\Gamma'$ for the basis
$a=srsr^{-1}$, $b=sr^2sr^{-2}$.
For a candidate conjugator $\beta$, fixing $sr$ makes the
image of $r$ equal to $dr$, where
$d=\beta\theta(\beta)^{-1}$.  Its cube is one exactly when
$d\alpha(d)\alpha^2(d)=1$.

\begin{proposition}[A truncated peripheral counterexample]
\label{prop:truncated-counterexample}
In the mod-two Magnus ring
$\mathcal R=\mathbb F_2\langle A,B\rangle/(\text{degree}\ge5)$,
send $a$ to $1+A$ and $b$ to $1+B$.
The word $\beta=a^3bab^3$ has exponent sums $(4,4)$ and
satisfies
\[
 d\alpha(d)\alpha^2(d)=1\quad\text{in }\mathcal R.
\]
Nevertheless its image under $a\mapsto i,b\mapsto j$ in the
quaternion group is $-1$.  Thus this truncated norm equation
and the two exponent sums modulo four do not determine the
quaternion central sign.
\end{proposition}
\begin{proof}
The exponent sums are immediate, and
$i^3jij^3=(ij)^2=-1$.
Here are direct expansions verifying the norm equation.  Write
$d=1+D$, $\alpha(d)=1+E$, $\alpha^2(d)=1+F$ in $\mathcal R$.
Substitute $a^{-1}=1+A+A^2+A^3+A^4$ and the analogous
expression for $b^{-1}$.  One obtains
\[
\begin{split}
D={}&AAB+ABB+BAA+BBA\\
 &+AAAB+AABA+AABB+ABAA+ABBA+ABBB\\
 &+BAAA+BAAB+BABB+BBAA+BBAB+BBBA,\\
E={}&ABB+BBA+AABB+ABAB+ABBB+BAAB\\
 &+BABA+BABB+BBAB+BBBA,\\
F={}&AAB+BAA+AAAB+AABA+ABAA+ABAB\\
 &+ABBA+BAAA+BABA+BBAA.
\end{split}
\]
Every monomial in $D+E+F$ occurs twice, so the sum is zero.
Each of $D,E,F$ has degree at least three, hence all their
pairwise products vanish in $\mathcal R$.  Therefore
$(1+D)(1+E)(1+F)=1$.
The trivial word has the same exponent sums modulo four
and passes the same norm equation, but has quaternion image
$+1$.  This proves the claimed limitation.
\end{proof}

The displayed expansion proves only the limitation of this truncated
norm equation.  It verifies neither the full pentagon nor the full
nonabelian distribution relation, and does not produce a profinite
GT element.

\subsection{Two-dimensional specializations of relation $(\mathrm{III}^\prime)$}

Write
\[
 B_4=\langle a,b,c\mid aba=bab,\ bcb=cbc,\ ac=ca\rangle.
\]
In the notation of \cite[\S10, Theorem~20]{H17}, the two inputs
to the harmonic auxiliary word $g_g=\beta_g^{-1}$ are
\[
 x_{45}=(ab)^3,\qquad x_{51}=(bc)^3.
\]
On the subgroup $D$, $g_g\in\widehat F(x,y)'$.

\begin{lemma}\label{lem:iii-prime-two-dimensional}
Let $k$ be any field and let $\rho:B_4\to\mathrm{GL}_2(k)$ be
a homomorphism.  Then $\rho((ab)^3)$ and $\rho((bc)^3)$ commute.
For a finite field $k$ and $g\in D$, it follows that
\[
 g_g\bigl(\rho((ab)^3),\rho((bc)^3)\bigr)=1.
\]
Thus such a direct matrix specialization of $(\mathrm{III}')$
does not retain its harmonic auxiliary-word term.
\end{lemma}
\begin{proof}
The elementary braid identity says that $(ab)^3$ is central in
$\langle a,b\rangle$, and $(bc)^3$ is central in
$\langle b,c\rangle$.  Both displayed images therefore belong
to the centralizer of $B=\rho(b)$.

If $B$ is not scalar, its centralizer in $M_2(k)$ is the
commutative algebra $k[B]$.  To see this directly, choose a vector
$v$ such that $v,Bv$ are a basis.  Such a vector exists for every
nonscalar two-by-two matrix.  An operator commuting with $B$ is
determined by its value on $v$, and the operators $\alpha I+\beta B$
give every possible such value.  In particular the two centralizer
elements commute.

If $B=tI$, the braid relation with $A=\rho(a)$ gives
$tA^2=t^2A$, and invertibility implies $A=tI$.  The other braid
relation similarly gives $\rho(c)=tI$.  All images then commute.

Over a finite field the representation extends continuously to
the profinite completion.  Evaluation of a commutator proword on
two commuting elements is one.  Applying this to $g_g$ proves the
second assertion.
\end{proof}

This lemma only excludes this direct way of retaining the
$g_g$ term.  It does not assert that $(\mathrm{III}')$ has no
consequences for the finite defects, nor that higher-dimensional
representations or comparisons of covering groups cannot provide
the required equality.  No vanishing or nonvanishing statement for
the two defects follows from it.


\subsection{A genus-two quotient carrying the cubic comparison sign}

Use the finite-comparison notation
$U=\left(\begin{smallmatrix}1&1\\0&1\end{smallmatrix}\right)$,
$L=\left(\begin{smallmatrix}1&0\\-1&1\end{smallmatrix}\right)$,
$R=-UL$, $X=U^2$, $Y=L^2$, $C=5I$, and $Z=-I$ modulo eight.
Write $G_8=\mathrm{SL}_2(\mathbb Z/8)$ and
$P=\langle X,R\rangle$, of order $192$.
The already proved presentation has
\[
 N=\langle X,Y,C,Z\rangle,\quad |N|=64,\quad
 P=N\rtimes\langle R\rangle,
\]
with the unique forms $X^aY^bC^eZ^d$,
$0\leq a,b<4$, $0\leq e,d<2$, and
\[
 [X,Y]=C,\qquad RXR^{-1}=Y,\qquad
 RYR^{-1}=Z(XY)^{-1}.
\]
The elements $C,Z,X^2,Y^2$ are central in $N$.

In the target put
\[
 S=\begin{pmatrix}0&1\\-1&0\end{pmatrix},\quad
 T=\begin{pmatrix}1&1\\0&1\end{pmatrix},\quad
 A=I-S=\begin{pmatrix}1&2\\1&1\end{pmatrix}
 \quad\text{over }\mathbb F_3.
\]

\begin{proposition}\label{prop:genus-two-finite-quotient}
There is a surjection
\[
 \ell:G_8\longrightarrow\mathrm{GL}_2(\mathbb F_3),\qquad
 U\longmapsto A,\quad R\longmapsto T,
\]
whose kernel is the order-eight subgroup
\[
 K=\langle Z,X^2C,Y^2C\rangle.
\]
It kills $Z$, so the signed congruence representation induces
\[
 \bar\ell:\Gamma=C_2*C_3\longrightarrow
                 \mathrm{GL}_2(\mathbb F_3).
\]
For $s=uvu$, $r=uv$ and $u=r^{-1}s$, their images have orders
$2,3,8$, respectively.  On the subgroup $P$, this map agrees
with the previous interchange map on all commutator prowords in
$X,R$, although its value on $X$ is $S$ instead of $-S$.
\end{proposition}
\begin{proof}
The presentation of $P$ defines a character
$\chi:P\to\{\pm1\}$ by
\[
 \chi(X)=\chi(Y)=\chi(Z)=-1,\qquad
 \chi(C)=\chi(R)=1.
\]
For example the only nontrivial conjugation check is
$\chi(Y)=\chi(Z(XY)^{-1})=-1$; the remaining relations are
immediate.  Multiply the previously established interchange
homomorphism $\kappa:P\to\mathrm{SL}_2(\mathbb F_3)$ by this
central character.  The resulting surjection $\kappa'$ has
\[
 \kappa'(X)=S,\quad \kappa'(R)=T,\quad
 \kappa'(C)=-I,\quad\kappa'(Z)=I.
\]
Its kernel contains $K$.  The unique normal forms show that
$Z,X^2C,Y^2C$ are three independent commuting involutions, so
$|K|=8$.  Surjectivity and $|P|=192$ imply that this is the whole
kernel.  A character is trivial on commutators, giving the last
assertion about prowords.

The subgroup $P$ has index two in $G_8$, and $U^2=X\in P$.
The whole group is generated by $P$ and $U$: indeed
$L=U^{-1}ZR$.  Direct multiplication gives
\begin{equation}\label{eq:bolza-1}
UXU^{-1}=X,\qquad URU^{-1}=XYR^2.
\end{equation}
In the target, $A^2=S$, $A$ commutes with $S$, and
\begin{equation}\label{eq:bolza-2}
ATA^{-1}=S(TST^{-1})T^2.
\end{equation}
For clarity, both sides of \eqref{eq:bolza-2} are
$\left(\begin{smallmatrix}2&2\\1&0\end{smallmatrix}\right)$.
Since $X,R$ generate $P$, \eqref{eq:bolza-1}--\eqref{eq:bolza-2} show that conjugation by
$A$ on $\kappa'(P)$ agrees with conjugation by $U$ before
applying $\kappa'$.  Together with $A^2=\kappa'(U^2)$, these
identities give a well-defined homomorphism on the two cosets
$P\sqcup PU$, extending $\kappa'$ and sending $U$ to $A$.
Its image contains $\mathrm{SL}_2(\mathbb F_3)$ and the matrix
$A$ of determinant $-1$, hence is all of $\mathrm{GL}_2(\mathbb F_3)$.
Its kernel has order $384/48=8$ and is $K$.

The signed braid center maps to $Z$, so the map descends to
$\Gamma$.  There $r=uv$ maps to $T$, since $ZR$ maps to $T$;
and $s=ru$ maps to
\[
 TA=\begin{pmatrix}2&0\\1&1\end{pmatrix},
\]
an involution.  Finally $A^2=S$, $S^2=-I$ and $A^4=-I$,
so $A$ has order eight.  This proves the stated orders.
\end{proof}

\begin{corollary}\label{cor:genus-two-sign-cover}
The discrete kernel of $\bar\ell$ is the fundamental group of
a genus-two surface with six punctures.  Its compactification
carries a faithful deck action of $\mathrm{GL}_2(\mathbb F_3)$
of signature $(2,3,8)$.  The central involution is its
hyperelliptic involution.

On $D$ the induced action on this finite deck group is the
identity.  The pure harmonic conjugator maps to $I$, and the
pure cubic conjugator maps to $(-I)^{\tau(g)}$.
\end{corollary}
\begin{proof}
The subgroup has index $48$ in the modular group $\Gamma$.
It is torsion-free because the orders two and three of the
finite factors are preserved.  The cusp stabilizer has image
of order eight, so there are $48/8=6$ cusps.  Since
$\chi(\Gamma)=-1/6$, its Euler characteristic is $-8$.
Thus $2-2g-6=-8$ and $g=2$.  Filling the cusps gives the
signature $(2,3,8)$ action, by the same stabilizer counts.

The central matrix $-I=A^4$ lies in every conjugate of the
cusp stabilizer, so it fixes all six filled cusp points.
For an involution of a genus-two compact surface the
Riemann--Hurwitz formula is $2=4h-4+r$.  Since $r\geq6$,
one has $h=0$ and $r=6$.  This is the hyperelliptic involution.

On $D$ the established normalized signed covariance gives
$q_8\Phi_g=q_8$, hence identity on the deck group after $\ell$.
The established prime-two normalization gives
$q_8(\beta_g)=I$.  Also
$q_8(A_g)^{-1}=5^{\tau(g)}I=C^{\tau(g)}$, and
$\ell(C)=-I$.  Inversion leaves this sign unchanged, yielding
the asserted cubic conjugator value.
\end{proof}

This gives a concrete genus-two setting for the cubic finite
comparison. The equality $\tau=\eta$ requires the cubic branch-filling
comparison proved in Section~\ref{sec:normalized-cubic-identity}.

\paragraph{Identification with the classical Bolza cover.}
The genus-two cover in Proposition~\ref{prop:genus-two-finite-quotient}
is the classical Bolza cover.  Katz--Katz--Schein--Vishne
\cite[\S3 and Corollaries 13.3--13.4]{KKSV2016} describe its
index-two triangle-group inclusion and the corresponding
$\mathrm{SL}_2(\mathbb F_3)$ and $\mathrm{GL}_2(\mathbb F_3)$
quotients.  Their displayed matrices immediately before
Corollary 13.4 send the order-two and order-three generators to
\[
 s\longmapsto\begin{pmatrix}0&1\\1&0\end{pmatrix},\qquad
 r\longmapsto T=\begin{pmatrix}1&1\\0&1\end{pmatrix}.
\]
These are exactly conjugate to our pair:
\[
 T(TA)T^{-1}=\begin{pmatrix}0&1\\1&0\end{pmatrix},\qquad
 TTT^{-1}=T.
\]
Thus the kernels in the triangle group agree.  The compactification is the Bolza
curve $y^2=x^5-x$; the same model, its six hyperelliptic branch
points $0,\infty,\pm1,\pm i$, and its classical automorphism
group are recalled in \cite[\S4, equation (4.1)]{KRRBolza}.
That paper attributes the classification to Bolza's 1887 work.

Here we compute the factorization through
$\mathrm{SL}_2(\mathbb Z/8)$ with kernel
$\langle Z,X^2C,Y^2C\rangle$, together with the values of the
normalized frame conjugators. Section~\ref{sec:normalized-cubic-identity}
then identifies their cubic sign with $\eta$.

